\chapter{Estado del arte}\label{cap5}

A continuación, se presentan algunos de los trabajos más recientes que se
encuentran en el campo de estudio de interés. Estos permitirán conocer el estado
actual de la investigación y desarrollo en el área de los biorreactores, así como
posicionar el presente trabajo en un contexto actual de los avances ya logrados.

En 2025, Mauro Moresi, catedrático titular en Operaciones de Ingeniería de
Alimentos y Evaluación del Impacto Ambiental en la Industria Alimentaria en la
Universidad de Tuscia, Italia, realizó un proyecto de investigación que culminó con
el desarrollo de un biorreactor de torre con camisa versátil. Este diseño permitió
optimizar la eficiencia en la transferencia de oxígeno y el control del proceso. Su
flexibilidad operativa permitió el funcionamiento en dos modos, el modo airlift
(circulación por aire) y en modo bucle eyector, ideal en varios de los procesos de
fermentación aeróbica. Al operar en modo airlift, se identificaron problemas debido
a una sedimentación no homogénea, lo que demuestra la necesidad de un sistema
de recirculación controlada. En modo bucle eyector resultó eficaz para algunas
cepas, pero inadecuado en organismos sensibles, puesto que los esfuerzos
cortantes resultaron perjudiciales \cite{Moresi2025}.

Este trabajo es de interés puesto que muestra un diseño versátil aplicable al cultivo
de organismos aerobios, como lo son los ascomicetos coprófilos. A lo largo del
mismo se resaltan las consideraciones de diseño ingenieril y las mejoras
propuestas, como la necesidad de recirculación controlada en futuros trabajos.

El trabajo realizado en 2025 por Federico Cerrone, profesor asistente de Ingeniería
de Bioprocesos en la Universidad de la Ciudad de Dublín e investigador principal
asociado en BiOrbic, y Kevin E O'Connor, catedrático titular en Bioeconomía y
Bioprocesos en la Universidad de la Ciudad de Dublín, Irlanda, llevaron a cabo un
estudio sobre el cultivo de hongos filamentosos en biorreactores de tipo airlift. El
trabajo analizó las ventajas y desventajas del cultivo de hongos filamentosos en
biorreactores de tipo airlift. En el estudio, se hace una comparativa de estos
biorreactores con otros medios de cultivo, como lo son los biorreactores de tanque
agitado y de columna de burbujeo.

El análisis concluye que el cultivo en biorreactores tipo airlift ofrece una relación
entre transferencia volumétrica de oxígeno y estrés por cizalladura, lo que aumenta
la productividad en comparación con los otros medios estudiados. Sin embargo, se
hace énfasis en la necesidad de un control adecuado del tamaño y densidad de los
agregados miceliales formados, ya que pueden disminuir la eficiencia del proceso.
En el artículo se mencionan también mejoras en el diseño para aumentar la
producción de hongos \cite{Cerrone2025}.

En 2023, Francisco Gabriel Pardo Pérez, Ingeniero Civil Químico de la Universidad
de Chile, asesorado por María Elena Lienqueo Contreras, PhD en Ingeniería
Química de la Universidad de Chile, presentó su investigación sobre una estrategia
de escalamiento de un biorreactor tanque agitado, con el fin de producir biomasa y
proteína enfocado de hongos filamentosos. Para lograrlo, se optimizó el control de
la temperatura, aireación y agitación; además de obtenerse el coeficiente
volumétrico de transferencia de masa ($k_L a$), con el objetivo de escalar sistema de
laboratorio a una planta industrial.

Este proyecto aporta un gran valor al desarrollo del trabajo propuesto, debido a que
muestra la metodología utilizada para escalar un proceso de fermentación
sumergida en laboratorio a uno industrial considerando variables como la
viscosidad del medio y esfuerzos de corte \cite{Pardo2023}.

En 2024, Federico Cerrone, investigador principal en el BiOrbic Bioeconomy
Research Centre y profesor en la University College Dublin, junto con Kevin E.\
O'Connor, catedrático en Bioeconomía y Bioprocesos, presentaron una
investigación en la que se analizan las estrategias de cultivo en biorreactores tipo
airlift aplicadas al crecimiento de Agaricomycetes, hongos filamentosos comestibles \cite{Cerrone2024}.

Este artículo documenta el diseño, operación y rendimiento de un biorreactor airlift,
de igual manera aborda aspectos clave al tratar con hongos filamentosos, como la
formación de pellets, consumo de oxígeno y acumulación de metabolitos. Esta
información será de utilidad ya que, a pesar de centrarse en Basidiomycetes y no
en Ascomicetos, las condiciones de cultivo en entornos controlados son similares.

En 2024, Diego Amador Guerra, Jorge Luis Chimal Ayala y Leví Alberto González
López, estudiantes de Ingeniería en Biotecnología del Instituto Tecnológico y de
Estudios Superiores de Occidente (ITESO), diseñaron el bioproceso completo para
la producción de extracto antioxidante presente en \textit{Lepista sordida}. En este trabajo
se seleccionó un reactor de tanque agitado, comparándolo con uno de tipo airlift,
destacando ventajas y desventajas. Este trabajo es útil porque presenta una
experiencia aplicada en la comparación de diseño de biorreactores, además de
presentar un marco teórico solido para la selección de medios de cultivo,
escalamiento y evaluación de productos bioactivos \cite{Amador2024}.
